University course timetabling must place activities in time and rooms while preserving student, staff, equipment, capacity, and institutional constraints.
The optimization problem is mature, but the operational failure mode remains difficult: a time-feasible schedule can still be impossible to room, and a bare ``infeasible'' status does not tell an administrator which activities must move.

Time/room decomposition is not new.
Bagger et al.\ explicitly decompose curriculum-based course timetabling into time scheduling and room allocation~\cite{bagger2018benders}; hybrid and metaheuristic approaches are also well established~\cite{muller2009hybrid,lewis2008metaheuristics,bettinelli2015overview}.
Fair timetabling based on max--min orderings and Jain's index has likewise been studied~\cite{muhlenthaler2013fairness,muhlenthaler2013decomposition}.
The research question in this work is therefore not whether decomposition or fairness works in general.
It is whether a room-subproblem certificate can serve three roles in one auditable system: a valid master cut, a concise explanation, and a bounded repair neighborhood.

The implemented contributions are:
\begin{itemize}[leftmargin=*]
  \item a CP-SAT time master coupled to an exact, full-domain room assignment subproblem;
  \item Hall-deficiency extraction for simultaneous room jobs, with capacity, room type, specialization, availability, closure, and lock filtering performed before the certificate is formed;
  \item iterative certificate cuts that remove a deficient occupancy pattern, plus exact nogoods for infeasibility caused by travel, repeat-room, or lock interactions not captured by a slot-wise Hall witness;
  \item reuse of the same certificate for human-readable explanations and a resource-linked repair neighborhood while freezing unaffected activities;
  \item an adaptive large-instance path that validates global invariants, solves decision-independent weeks concurrently, accepts only validated room assignments, activates exact certificate decomposition on a failing week, and refuses hard cross-week coupling rather than silently weakening it;
  \item enrollment-uncertainty policies whose nominal, scenario-quantile, worst-case, or budgeted demand is reused consistently by room domains, validators, metrics, and infeasibility certificates;
  \item a portable distribution-constraint schema with exact evaluation for nineteen common ITC-style relations, hard CP compilation for seventeen of them, and fail-fast handling instead of silent weakening for the two evaluator-only aggregate rules;
  \item executable institution-policy presets, including a portable baseline and a repository-derived GIU target preset that is explicitly marked as requiring local policy validation;
  \item an exact lexicographic fairness profile that minimizes the maximum group burden before total weighted penalty, accompanied by tail, Gini, and Jain diagnostics;
  \item a reproducible evaluation harness with effective instance/solver/local-search seeds, instance and schedule fingerprints, solver bounds and gaps, paired statistics, official ITC-2007 soft scoring, ITC-2019 schema inspection, portable metrics, disruption experiments, and verified release bundles.
\end{itemize}

The systematic review falsifies the broad mathematical novelty claim: the contextual cut is an option-expanded instance of the established partial-transversal Hall system, while Benders-guided and explanation-guided LNS both have direct prior art.
The remaining research hypothesis is a systems and empirical one: whether exact effective-domain reconstruction, conservative cut fallback, typed certificate-support transfer, reusable exact repair, and auditable lineage improve performance and operational trust when implemented together.
A publication submission must still complete the evaluation protocol in Section~VII and an independent literature review.
This paper deliberately separates implemented capability, verified tests, planned experiments, and unsupported claims.
